% titel: Blatt 10, A3: Säuren/Basen klassifizieren
% Teil II: der durchdachte Weg.

\subsection*{Einordnung}
\begin{itemize}
  \item \textbf{Na$_2$CO$_3$:} Brønsted-Base (CO$_3^{2-}$ nimmt H$^+$ auf);
        als Anion auch Lewis-Base (freies Paar an O).
  \item \textbf{AlCl$_3$:} Lewis-Säure (Elektronenpaar-Akzeptor, kein H);
        keine Brønsted-Säure.
  \item \textbf{Kalium-tert-butanolat:} gibt tBuO$^-$; Brønsted-Base
        (starke Base, nimmt H$^+$ auf) und Lewis-Base.
  \item \textbf{Aceton:} Lewis-Base (freies Paar an O); Brønsted-Base nur
        sehr schwach.
  \item \textbf{HCN:} Brønsted-Säure (gibt H$^+$); CN$^-$ bzw.\ N-Freies Paar
        macht es auch zur Lewis-Base.
  \item \textbf{BF$_3$:} Lewis-Säure (Elektronenpaar-Akzeptor), kein H.
  \item \textbf{Diethylether:} Lewis-Base (freies Paar an O); auch schwache
        Brønsted-Base.
  \item \textbf{HNO$_3$:} starke Brønsted-Säure (gibt H$^+$); die
        Nitrat-Anionen sind auch Lewis-Basen.
\end{itemize}

\begin{ergebnis}
Brønsted-Säuren: HCN, HNO$_3$ (und je nach Kontext partielle).
Lewis-Säuren: AlCl$_3$, BF$_3$.
Basen (Brønsted+Lewis): Na$_2$CO$_3$, tert-Butanolat, Aceton, Diethylether,
z.\,T.\ HCN/HNO$_3$ (als Anionen).
\end{ergebnis}

\begin{kontrolle}
Die Zuordnung folgt sauber der Definition: fehlt H (BF$_3$, AlCl$_3$)
$\Rightarrow$ reine Lewis-Säure; freies Paar $\Rightarrow$ Lewis-Base;
H-Donator $\Rightarrow$ Brønsted-Säure. Kein Widerspruch.
\end{kontrolle}

\begin{merksatz}
Brønsted verlangt H$^+$-Transfer, Lewis nur Elektronenpaar-Vermittlung.
Lewis-Säuren müssen kein H enthalten; Anionen und O/N-haltige Moleküle sind
Lewis-Basen.
\end{merksatz}
